\documentclass[11pt]{article}
\usepackage[a4paper,margin=30mm]{geometry}
\usepackage[T1]{fontenc}
\usepackage{lmodern}
\usepackage{microtype}
\usepackage{amsmath,amssymb,amsthm,mathtools}
\usepackage{booktabs}
\usepackage{hyperref}
\usepackage{xcolor}
\hypersetup{
  colorlinks=true,
  linkcolor=blue!45!black,
  citecolor=blue!45!black,
  urlcolor=blue!55!black,
  pdftitle={Exact Finite Rigidity for Furter's R(3): All Windows Through n=299},
  pdfauthor={Anonymous},
  pdfsubject={An exact computer-assisted theorem for Furter's R(3,n), 1 <= n <= 299},
  pdfkeywords={polynomial composition rigidity, Polydegree Conjecture, Strong Factorial Conjecture, computer-assisted proof}
}

\newtheorem{theorem}{Theorem}[section]
\newtheorem{proposition}[theorem]{Proposition}
\newtheorem{lemma}[theorem]{Lemma}
\newtheorem{corollary}[theorem]{Corollary}
\newtheorem{remark}[theorem]{Remark}
\theoremstyle{definition}
\newtheorem{definition}[theorem]{Definition}
\newcommand{\QQ}{\mathbb{Q}}
\newcommand{\ZZ}{\mathbb{Z}}
\newcommand{\FF}{\mathbb{F}}
\newcommand{\GG}{\mathcal{G}}

\title{Exact Finite Rigidity for Furter's \texorpdfstring{$R(3)$}{R(3)}:\\All Windows Through \texorpdfstring{$n=299$}{n=299}}
\author{Anonymous}
\date{14 August 2026\\Version 0.1.0\\\href{https://doi.org/10.5281/zenodo.21939362}{doi:10.5281/zenodo.21939362}}

\begin{document}
\maketitle

\begin{abstract}
Let $g_{k,3}\in\ZZ[x_1,x_2,x_3]$ be the coefficient of $X^{k+1}$ in the
compositional inverse of
$X(1+x_1X+x_2X^2+x_3X^3)$.  We prove, by exact computer-assisted
certificates, that
\[
 \sqrt{(g_{d,3},g_{d+1,3},g_{d+2,3})\QQ[x_1,x_2,x_3]}
   =(x_1,x_2,x_3)
 \qquad(2\le d\le300).
\]
Equivalently, Furter's polynomial-composition rigidity assertion $R(3,n)$
holds for every $1\le n\le299$.  This extends the previously published
finite computation for the relevant three-variable window from $d<50$ to
$d\le300$.  Each new instance is certified at one prime by exact unit-ideal
computations on all three strata of a weighted-projective cover; properness
then transfers emptiness of the special fibre to characteristic zero.  We
also prove a periodic Kummer-theoretic zero-band theorem showing that no
single prime can certify every degree by this route.  Exact rational
Euler--socle certificates through $d=12$ and normalized-resultant data through
$d=25$ isolate two structural continuations, but universal $R(3)$ remains
open.  A frozen package supplies ordinary and optimized replay, independent
coefficient reconstruction, negative controls, hashes, and a fresh-extraction
verification receipt.
\end{abstract}

\noindent\textbf{Status.} This is an unrefereed computer-assisted theorem
candidate.  The exact finite claim is stronger than a probabilistic
calculation, but the release has not received specialist peer review or an
independent external reconstruction.

\section{Introduction}

Furter's rigidity conjectures ask whether an unexpectedly long initial
cancellation in a polynomial composition forces both polynomials to be the
identity.  They are important because Furter connected them to closure
relations between strata of plane polynomial automorphisms
\cite{Furter2015}.  Edo and van den Essen found a further bridge to a
restricted family in the Strong Factorial Conjecture \cite{EdoVdE2014}.
The universal statement $R(3)$ is still open.

Lewis, Perry and Straub (LPS) recast related closure questions as explicit
ideal computations in inverse-series coefficients and verified their stronger
condition in the three-variable column for $d<50$ \cite{LPS2019}.  The present
work takes the radical-maximality half of that architecture and extends its
exact finite rigidity consequence to $d\le300$.

Our principal result is deliberately finite.

\begin{theorem}[Exact finite $R(3)$]\label{thm:main}
Furter's $R(3,n)$ holds for every integer $n$ with $1\le n\le299$.
\end{theorem}

The proof does not infer a characteristic-zero theorem from the apparent
frequency of good modular reductions.  For every degree it records an exact
empty special fibre over a single finite field.  A properness argument then
gives the corresponding characteristic-zero conclusion.  A bad prime has no
negative implication for characteristic zero; it merely cannot certify that
degree by this route.

The computation reveals a structural limitation.  For every prime $p\ge7$,
infinitely many inverse coefficients vanish identically modulo $p$, so
infinitely many three-term windows are necessarily bad at that prime.  The
exact map supplies a bad window for each of $p=2,3,5$ as well.  Thus no fixed
prime certifies every degree, and any universal proof based on good special
fibres must coordinate multiple primes.

\section{Inverse coefficients and the rigidity dictionary}

Give $x_i$ weight $i$.  Define $g_{n,3}$ by
\begin{equation}\label{eq:g-explicit}
g_{n,3}=\sum_{a_1+2a_2+3a_3=n}
(-1)^{a_1+a_2+a_3}
\frac{(n+a_1+a_2+a_3)!}
{(n+1)n!a_1!a_2!a_3!}
x_1^{a_1}x_2^{a_2}x_3^{a_3}.
\end{equation}
These are integral weighted-homogeneous polynomials of weight $n$.  Equivalently,
if $G$ is the unique formal solution with initial term $w$ of
\begin{equation}\label{eq:inverse-series}
G=w-x_1G^2-x_2G^3-x_3G^4,
\end{equation}
then $g_{n,3}$ is the coefficient of $w^{n+1}$ in $G$.

For $d\ge2$ put
\[
 I_d=(g_{d,3},g_{d+1,3},g_{d+2,3})\subset\ZZ[x_1,x_2,x_3].
\]

\begin{lemma}[Source dictionary]\label{lem:dictionary}
The equality
\[
 \sqrt{I_d\QQ[x_1,x_2,x_3]}=(x_1,x_2,x_3)
\]
is Furter's one-map assertion $\widetilde R(3,d-1)$ and is equivalent to his
two-map assertion $R(3,d-1)$.
\end{lemma}

\begin{proof}
The coefficient convention follows from Lagrange inversion and agrees with
the LPS definition.  Furter's Lemma~2 identifies the vanishing of the three
consecutive inverse coefficients with $\widetilde R(3,d-1)$ and proves the
equivalence with $R(3,d-1)$ \cite{Furter2015,LPS2019}.  The only common
affine zero being the origin is precisely the displayed radical equality.
\end{proof}

The shift $n=d-1$ is essential: $2\le d\le300$ corresponds exactly to
$1\le n\le299$.

\section{From one good special fibre to characteristic zero}

Let $p$ be prime and write $I_{d,p}$ for the reduction of $I_d$ in
$\FF_p[x_1,x_2,x_3]$.

\begin{definition}
The pair $(d,p)$ is \emph{GOOD} if $I_{d,p}$ has height three, equivalently
if the quotient is finite-dimensional.  It is \emph{BAD} if the three forms
have a nonzero common point over $\overline{\FF}_p$.
\end{definition}

Because the forms have weighted degrees $d,d+1,d+2$, a GOOD quotient is a
weighted complete intersection.  Its length is
\begin{equation}\label{eq:length}
 \frac{d(d+1)(d+2)}{1\cdot2\cdot3}=\binom{d+2}{3}
\end{equation}
and its socle weight is
\begin{equation}\label{eq:socle}
 d+(d+1)+(d+2)-(1+2+3)=3d-3.
\end{equation}

\begin{lemma}[Good-special-fibre bridge]\label{lem:properness}
If $(d,p)$ is GOOD, then
$\sqrt{I_d\QQ[x_1,x_2,x_3]}=(x_1,x_2,x_3)$.
\end{lemma}

\begin{proof}
Regard the common-zero scheme as a closed subscheme of the weighted projective
plane over $\operatorname{Spec}\ZZ_{(p)}$.  This morphism is proper.  If its
generic fibre were nonempty, the closure of a generic point would have closed
image containing the generic point of $\operatorname{Spec}\ZZ_{(p)}$, hence
the closed point as well.  The special fibre would then be nonempty, contrary
to GOOD.  Thus the generic weighted-projective fibre is empty, so the only
affine common zero in characteristic zero is the origin.
\end{proof}

This argument uses one prime for the whole projective fibre.  Certificates
from different primes on different affine charts cannot be combined.

For completeness, the same bridge has a graded-module formulation.  With
$A=\ZZ[x_1,x_2,x_3]/I_d$, every weighted piece $A_k$ is a finite
$\ZZ$-module.  A GOOD special fibre is a complete intersection with Hilbert
series
\[
 \frac{(1-t^d)(1-t^{d+1})(1-t^{d+2})}
 {(1-t)(1-t^2)(1-t^3)}.
\]
The special fibre vanishes above weight $3d-3$.  Nakayama's lemma then kills
the corresponding localized integral pieces and makes the generic fibre
Artinian.  Equal complete-intersection Hilbert functions also show that the
localized graded pieces are free, a fact used by the structural pilot to lift
modular nonvanishing.

\section{The exact finite computation}

Two proof-bearing lanes cover the degree interval.

\subsection{Complete modular map for $2\le d\le49$}

For every one of the 48 degrees and every one of the 168 primes $p\le997$,
an exact Singular Gr\"obner-basis computation classified the quotient as
GOOD or BAD.  The map contains 8,064 rows: 6,931 GOOD and 1,133 BAD.  An
independent formal-reversion implementation checks the generator formula on
a fixed grid.  The verifier rejects an impossible finite length and includes
a mutation in which the third generator is replaced by zero.

This lane reproduces the radical half of the published LPS range.  Its role in
the present release is regression and continuity; the new finite range begins
at $d=50$.

\subsection{Weighted-projective cover for $50\le d\le300$}

Materializing complete quotient bases becomes increasingly expensive.  The
second lane instead certifies the absence of a nonzero projective common zero
on the three strata
\begin{equation}\label{eq:charts}
 x_1=1,\qquad x_1=0,\ x_2=1,\qquad
 x_1=x_2=0,\ x_3=1.
\end{equation}
These strata cover the weighted projective plane over an algebraically closed
field, including in characteristics $2$ and $3$.  For a fixed attempt $(d,p)$,
the producer reduces the integral generators at $p$ and proves that the ideal
on each stratum is the unit ideal.  All three results are nested under the
same prime-labelled attempt.  A second specialize-first implementation
rebuilds every selected GOOD certificate under ordinary and optimized Python.

The screen covers 251 degrees and records 2,854 attempts: 2,539 exact BAD
fibres, 251 selected GOOD fibres, and 64 timeouts.  A timeout has no
mathematical meaning and is never used as evidence.  The largest selected
prime is $127$, at $d=280$; because some earlier attempts timed out, the
selected prime is not always claimed to be the smallest possible GOOD prime.

\begin{proposition}[Certified finite range]\label{prop:finite}
For every $d$ with $2\le d\le300$, there exists a recorded prime $p$ for
which $(d,p)$ is GOOD.
\end{proposition}

\begin{proof}[Computer-assisted proof]
For $d\le49$, the exact modular-map receipt contains at least one row labelled
GOOD for every degree and the verifier recomputes the expected length
\eqref{eq:length}.  For $50\le d\le300$, the assembled projective-screen
receipt contains exactly one selected same-prime three-stratum unit-ideal
certificate for every degree.  Independent ordinary and optimized verifier
receipts agree byte-for-byte.  A root verifier binds the degree inventory,
source and runtime hashes, controls, and subordinate receipts.  A package
verifier then extracts the frozen archive into a fresh directory, runs the
replay, checks the manifest, and reproduces a PASS root.  All computations are
exact; no floating-point tolerance enters an acceptance path.
\end{proof}

\begin{proof}[Proof of Theorem~\ref{thm:main}]
Apply Proposition~\ref{prop:finite} and Lemma~\ref{lem:properness} for every
$2\le d\le300$, then use Lemma~\ref{lem:dictionary} with $n=d-1$.
\end{proof}

\begin{corollary}[Length-two Polydegree consequence]\label{cor:polydegree}
For every $2\le k\le300$, Furter's closure equality of Theorem~B holds for
the length-two multidegrees $(4,k)$ and $(k,4)$.
\end{corollary}

\begin{proof}
Use $R(3,k-1)$ in Furter's Theorem~B and his symmetry lemma
\cite{Furter2015}.
\end{proof}

This corollary is not presented as the first proof of the full $e=3$
Polydegree column.  A separate Evidence Press candidate has already proved
that column for all degrees by a smooth-point/Fourier/interval-arithmetic
route \cite{FullE3Candidate}.

\section{A fixed-prime obstruction}

The modular data suggests blocks, but the key negative fact is analytic.

\begin{theorem}[Periodic zero band]\label{thm:zeroband}
Let $p\ge7$ be prime and put $K_p=\lfloor(p+2)/4\rfloor$.  If
\[
 n\equiv-k\pmod p,\qquad 2\le k\le K_p,
\]
then $g_{n,3}$ is the zero polynomial in
$\FF_p[x_1,x_2,x_3]$.
\end{theorem}

\begin{proof}
Consider a monomial indexed by $(a_1,a_2,a_3)$ in
\eqref{eq:g-explicit}, write $s=a_1+a_2+a_3$, and let $b_i$ be the units
digit of $a_i$ in base $p$.  The weighted-degree condition gives
\[
 b_1+2b_2+3b_3\equiv-k\pmod p.
\]
If $b_1+b_2+b_3<k$, then the left side is at most $3(k-1)<p-k$, a
contradiction.  Thus $b_1+b_2+b_3\ge k$.  Adding the base-$p$ units digits
of $n,a_1,a_2,a_3$ therefore produces a carry, because the units digit of
$n$ is $p-k$.  By Kummer's theorem the multinomial coefficient
\[
 \frac{(n+s)!}{n!a_1!a_2!a_3!}
\]
is divisible by $p$.  Since $n+1\equiv1-k\not\equiv0\pmod p$, division by
$n+1$ preserves divisibility.  Every coefficient of $g_{n,3}$ vanishes.
\end{proof}

\begin{corollary}[No universal fixed GOOD prime]\label{cor:nofixedprime}
For every prime $p\ge7$, infinitely many windows $I_{d,p}$ are BAD.  No
prime, including $2,3,5$, supplies a GOOD witness for every degree through
the good-special-fibre architecture.
\end{corollary}

\begin{proof}
For $p\ge7$, any three-term window intersecting the periodic zero band has at
most two nonzero generators and thus height at most two.  Explicitly,
\[
 p-K_p-2\le d\bmod p\le p-2
\]
forces BAD, and each congruence class recurs infinitely often.  For
$p=2,3,5$, the exact map supplies a BAD witness at $d=2,3,6$ respectively.
Those three witnesses are sufficient for the stated no-universal-fixed-prime
conclusion; no infinitude assertion for the small primes is needed here.
\end{proof}

The same carry proof works for weights $1,\ldots,e$.  If
$K_{p,e}=\lfloor(p+e-1)/(e+1)\rfloor$, then
$g_{n,e}=0\pmod p$ whenever $n\equiv-k\pmod p$ and
$2\le k\le K_{p,e}$.  We make no priority claim for this elementary
generalization.

\section{Structural evidence beyond the finite theorem}

\subsection{The canonical Euler--socle route}

Define the weighted Euler matrix
\[
 (M_d)_{ij}=\frac{j}{d+i-1}\partial_{x_j}g_{d+i-1,3}.
\]
Weighted Euler identities give
\[
 M_d(x_1,x_2,x_3)^T=(g_{d,3},g_{d+1,3},g_{d+2,3})^T.
\]
Consequently the adjugate identity places $x_i\det(M_d)$ in $I_d$ for each
$i$, while $\det(M_d)$ has the socle weight $3d-3$.  In an Artinian complete
intersection the Jacobian represents a nonzero socle class
\cite{SchejaStorch1975}.

This contracts one possible all-degree proof to two canonical identities.  It
would suffice to prove, for every $d\ge2$, nonzero scalars $\lambda_d,\nu_d$
such that
\begin{equation}\label{eq:euler-targets}
 x_1^{3d-3}-\lambda_d\det(M_d)\in I_d,
 \qquad
 x_3^{d-1}-\nu_d\det(M_d)\in I_d.
\end{equation}
Indeed the adjugate identities would force $x_1=x_3=0$ on the common-zero
locus, and the nonzero pure-$x_2$ Fuss--Catalan coefficient in the even member
of each three-term window would then force $x_2=0$.

Exact rational certificates verify the three pure-power memberships, the
one-lower nonmemberships, the Hilbert functions, the Euler identities, and
\eqref{eq:euler-targets} for $2\le d\le12$.  This is a finite structural
pilot, not an interpolation argument.

\subsection{The finite-prime automaton route}

Equation~\eqref{eq:inverse-series} defines an algebraic multivariate series
over every finite field.  Multivariate Christol theory therefore gives a
finite $p$-kernel for its coefficient array \cite{AdamczewskiBostanCaruso,
RowlandYassawi}.  What is missing is not coefficient automaticity, but a
finite-state lift from coefficient sections to an ideal certificate, such as
a bounded family of leading-monomial or Gr\"obner signatures.  In view of
Corollary~\ref{cor:nofixedprime}, a successful construction must use a product
automaton for several primes and show that no reachable state is BAD in every
component.

\subsection{Calibrated negative results}

Normalized weighted resultants were independently reconstructed through
$d=25$: full integer determinants through $d=17$ and modular determinants
with rigorous CRT reconstruction for $18\le d\le25$.  Across 4,032
prime-degree comparisons, a prime divides the normalized resultant exactly
when the corresponding fibre is BAD.  A withheld-data screen excludes every
first-order rational-ratio recurrence
$R_{d+1}Q(d)=R_dP(d)$ with $\deg P+\deg Q\le12$.  Separate screens reject a
small $d$-independent raw-cofactor transport and low-complexity recurrences for
the Euler--socle scalars.  These results close only the specified ansatz
classes; they do not rule out higher-order creative telescoping or a finite
automaton.

\section{Evidence object and assurance boundary}

The release archive is an executable evidence object rather than a table of
claimed CAS outputs.  Its root receipt binds:
\begin{itemize}
\item the 299-degree inventory and the two exact certificate lanes;
\item producer, verifier, input, runtime, and receipt SHA-256 identities;
\item ordinary/optimized Python parity where applicable;
\item independent formal reversion and sparse arithmetic on bounded pilots;
\item generator-axis Catalan and Fuss--Catalan regressions through $g_{302,3}$;
\item fail-closed controls for impossible colengths, missing generators,
  altered cofactors, incomplete prime prefixes, and malformed inventories;
\item an exact counterexample gate with no numerical-tolerance path; and
\item a fresh-extraction replay and manifest verification.
\end{itemize}

The large $d=50,\ldots,300$ lane independently replays a second Singular
wrapper, but it does not yet export portable sparse Nullstellensatz cofactors
for all 251 degrees.  Singular is therefore inside that lane's trusted
computing boundary.  Internal cross-model adversarial review is not external
specialist review, and executable replay is not formal verification.

\section{Limitations and open problem}

Theorem~\ref{thm:main} is finite.  It neither proves nor refutes the next
window, and it supplies no monotonicity principle in $d$.  Universal $R(3)$
would require an all-degree theorem, not a longer table.  The two best-defined
targets are:
\begin{enumerate}
\item prove both canonical identities \eqref{eq:euler-targets} for every
  $d\ge2$, including nonvanishing of their scalars; or
\item construct a finite set of primes and a certified product automaton that
  gives a same-prime GOOD fibre in every degree despite each prime's zero
  bands.
\end{enumerate}
The companion Universal $R(3)$ Challenge records these obligations, exact
finite examples, falsifiers, and acceptance criteria in machine-readable
form.

\section*{Data and code availability}

All proof scripts, receipts, controls, manifests, the frozen replay archive,
and the machine-readable challenge are included with the public release.  The
archive is identified by SHA-256 in its manifest and external verification
receipt.  The versioned source and replay package are available at
\url{https://github.com/ipitchford/furter-r3-through-299}; the permanent
release is \url{https://doi.org/10.5281/zenodo.21939362}.

\section*{Ethics statement}

This work uses no human participants, personal data, animals, clinical data,
or field interventions.  No ethics approval was required.

\section*{CRediT contribution statement}

Anonymous: conceptualization, methodology, software, validation, formal
analysis, investigation, data curation, writing, visualization, and project
administration.  The provenance record separately describes human research
direction and model-assisted implementation and review.

\section*{Competing interests}

The scholarly creator declares no competing interests.  Evidence Press is the
publisher of this unrefereed candidate and of related computer-assisted
mathematics releases.

\section*{Funding}

No external funding is declared.

\section*{AI assistance disclosure}

AI systems contributed to literature triage, mathematical exploration,
software implementation, adversarial checking, exposition, and release
engineering under human direction.  Exact computations were accepted only
through the recorded deterministic checks and replay gates.  Model agreement
was not treated as mathematical evidence, independent reproduction, or peer
review.

\begin{thebibliography}{99}

\bibitem{Furter2015}
J.-P. Furter,
\emph{Polynomial composition rigidity and plane polynomial automorphisms},
J. London Math. Soc. 91 (2015), 180--202.
\url{https://doi.org/10.1112/jlms/jdu064}

\bibitem{LPS2019}
D. Lewis, K. Perry and A. Straub,
\emph{An algorithmic approach to the Polydegree Conjecture for plane
polynomial automorphisms},
J. Pure Appl. Algebra 223 (2019), 5346--5359.
\url{https://doi.org/10.1016/j.jpaa.2019.04.002}

\bibitem{EdoVdE2014}
E. Edo and A. van den Essen,
\emph{The Strong Factorial Conjecture},
J. Algebra 397 (2014), 443--456.
\url{https://doi.org/10.1016/j.jalgebra.2013.09.011}

\bibitem{SchejaStorch1975}
G. Scheja and U. Storch,
\emph{\"Uber Spurfunktionen bei vollst\"andigen Durchschnitten},
J. Reine Angew. Math. 278/279 (1975), 174--190.
\url{https://eudml.org/doc/151652}

\bibitem{AdamczewskiBostanCaruso}
B. Adamczewski, A. Bostan and X. Caruso,
\emph{A sharper multivariate Christol's theorem with applications to
diagonals and Hadamard products}.
\url{https://arxiv.org/abs/2306.02640}

\bibitem{RowlandYassawi}
E. Rowland and R. Yassawi,
\emph{Automatic congruences for diagonals of rational functions},
J. Th\'eor. Nombres Bordeaux 27 (2015), 245--288.
\url{https://arxiv.org/abs/1310.8635}

\bibitem{FullE3Candidate}
Anonymous,
\emph{A Jacobian smooth-point criterion and the full $(e=3)$ column of the
Polydegree Conjecture}, unrefereed candidate, Evidence Press (2026).
\url{https://doi.org/10.5281/zenodo.21909085}

\end{thebibliography}

\end{document}
